\documentclass[11pt]{article}
\usepackage[utf8]{inputenc}
\usepackage[T1]{fontenc}
\usepackage[final]{pdfpages} 
\usepackage[french]{babel}
\usepackage{amsmath}
\usepackage[bookmarks={true},bookmarksopen={true}]{hyperref}
\usepackage{graphicx}
\usepackage[a4paper]{geometry}
\usepackage{listings}
	\lstset{frame=tb,
		language=Java,
 		aboveskip=3mm,
  		belowskip=3mm,
  		showstringspaces=false,
  		columns=flexible,
  		basicstyle={\small\ttfamily},
  		numbers=none,
 		numberstyle=\tiny\color{gray},
  		keywordstyle=\color{blue},
  		commentstyle=\color{dkgreen},
  		stringstyle=\color{mauve},
  		breaklines=true,
  		breakatwhitespace=true
  		tabsize=3
	}
\pagestyle{plain}
\setlength{\parindent}{5mm}

\usepackage{color}

\definecolor{dkgreen}{rgb}{0,0.6,0}
\definecolor{gray}{rgb}{0.5,0.5,0.5}
\definecolor{mauve}{rgb}{0.58,0,0.82}



\title{\textbf{Projet LSINF1121 -  Algorithmique et structures de données\\ - \\ Rapport final Mission 5} \\ {\large Groupe 3.2}}
\author{Boris \bsc{Dehem} \\(5586-12-00)\and Sundeep \bsc{Dhillon} \\(6401-11-00)\and Alexandre \bsc{Hauet} \\ (5336-08-00) \and Jonathan \bsc{Powell}\\(6133-12-00)\and Mathieu \bsc{Rosar} \\ (4718-12-00)\and Tanguy \bsc{Vaessen} \\ (0810-14-00)}
\date{date}
\date{\vspace*{25mm}
\includegraphics[scale=0.75]{logo.jpg}\\
		\vspace*{30mm}
		\begin{center}
		Année académique 2014-2015 \\	
		\end{center}}

\begin{document}
\thispagestyle{empty}

\maketitle
\thispagestyle{empty}
%\tableofcontents
\setcounter{tocdepth}{3}

\newpage
\setcounter{page}{1}

\section*{Introduction}

Dans le cadre du cours "Algorithmique et structures de données", il nous a été demandé de concevoir et d'implémenter une application permettant de compresser et décompresser des fichier textuelles sans perte d'information suivant l'algorithme de Huffman. Ainsi un code de longueur variable est attribué à chaque lettre du texte permettant d'avoir un codage en moyenne plus court que le fichier original. Ce programme contient une classe \verb+Compress.Java+ et \verb+Decompress.java+ qui prennent deux arguments un fichier à compresser/à décompresser et en deuxième argument le nom du fichier résultat.

\section{Implémentation}
Pour plus de facilité nous avons choisi que notre programme contiendrait deux \verb+Main+ différentes, une pour la compression et l'autre pour la décompression. Le programme Compress prendra en premier argument un fichier de type TXT et retournera un fichier compressé de format HUF.
Inversement, Decompress prendra un fichier compressé de format HUF et retournera un fichier TXT.

La classe \verb+HuffmanBTRee+ nous permet de créer un arbre selon l'algorithme de Huffman avec une queue de priorité implémentée par la classe \verb+PriorityQueue+ provenant de l'API Java.

Pour la compression, on parcoure le fichier lettre par lettre et pour chaque lettre rencontrée on la map dans une \verb+HasMap+ associé à sa fréquence qui sera incrémentée à chaque fois que la lettre est de nouveau rencontrée.

Ensuite, un \verb+HuffmanTree+ est créé par lettre associé à sa fréquence et stocké en fonction de la valeur de celle-ci dans une file de priorité.

Grâce à la fonction mergeAll de la classe \verb+HuffmanBTRee+ on va créer notre arbre final de Hufman correspondant au fichier qu'on souhaite compressé.

Cet arbre sera mis sous forme binaire, chaque nœud correspondra à un 0 et chaque feuille à un 1 suivi de la valeur du caractère en bits. Il sera ensuite incorporé au début du fichier compressé avec la taille du fichier original afin de pouvoir le décompresser plus tard.

Le fichier compressé aura la structure suivante : \\
$\left. \text{\parbox{0.76\linewidth}{
\begin{itemize}
\item <Integer> Nombre de caractères compressés écrits
\item <HuffmanBTree> Représentation binaire du huffmanBTree
\end{itemize}
}} \right \}$ Header 
\begin{itemize}
\item <Bits> Les chars compressés \\
\end{itemize} 

Pour la décompression on récupère le header avec les deux données qu'elle contient et on reconstruit un objet HuffmanBTree associer à ce fichier. Ensuite le fichier on le fichier compressé bit par bit, en même temps on parcourra l'arbre jusqu'à arriver à une feuille. Cette feuille correspondra à une lettre que sera écrit dans le fichier de sortie et la valeur du nombre de lettre du fichier orignal sera décrémenter de 1. Ainsi de suite à chaque lettre trouvé on recommence à parcourir l'arbre du début jusqu'à se que la valeurs totale de lettre e fichier orignal soit 0.


\section{Diagramme UML de classe}
Voir question 8 des questions liées au problème posé.

\section{Questions liées au problème posé}

\subsection*{Question 7}
En quoi les deux classes qui vous sont fournies, \textbf{InputBitStream} et OutputBitStream, peuvent-elles être utiles pour le problème de compression et de décompression avec un codage de Huffman ?
La postcondition de la méthode close dans la classe \textbf{OutputBitStream} précise notamment que \textit{si le nombre de bits déjà écrits ne correspond pas à un multiple de 8 (un octet), des bits à 0 sont écrits pour compléter l’octet courant}. Quand la situation décrite peut-elle se présenter ? Quelle est la conséquence de cette postcondition sur votre programme de compression de texte ? Quelle est la conséquence de cette postcondition sur votre programme de décompression ? \\


Une des particularités du codage de Huffman est qu'il n'utilise pas un nombre de bits constant pour encoder un caractère. Pour chaque fichier différent, l'arbre de Huffman le sera aussi. Grâce aux classes \verb+InputBitStream+ et \verb+OutputBitStream+ les caractères pourront être encodés et lu bit par bit sans avoir à connaître leur taille à l'avance.

L'encodage étant en \verb+ASCII+ ou \verb+Unicode+ représenté sur 16 bits, la partie du code contenant la représentation binaire du code ne sera pas influente sur ce point. En revanche, si jamais le nombre de bits utilisés pour représenter les différentes branches de l'arbre de Huffman dans le fichier compressé n' est pas un multiple de 8, alors la fermeture du stream nécéssitera de compléter l'octet courant entamé.

Cette postcondition est sans effet sur notre programme de compression, étant donné que la classe \verb+OutputBitStream+ avec sa méthode close() permet cette auto-complétion du dernier byte entamé, sans aucune ligne de code additionnelle de notre part.

Enfin, cette postcondition aura une influence sur notre programme de décompression, puisqu'il nous faudra alors faire attention à ne pas interpréter ces zéros supplémentaires comme des bits faisant partie de la représentation des caractères compressés. En résumé, il faut éviter l'erreur d'ajouter au fichier original des lettres supplémentaires issues de cette auto-complétion du dernier byte entamé.

\subsection*{Question 8}
Donnez un diagramme dé-taillé des différentes classes et interfaces utilisées pour résoudre le problème de compression et de décompression de données. Précisez dans ce diagramme, les liens du type \textit{uses}, \textit{implements} ou \textit{extends}. Quelles sont les classes qui sont communes entre les programmes de compression et de décompression de données ? \\

Le diagramme de classe se trouve en annexe (oir DiagMisson5.png). La seule classe commune est la classe \verb+HuffmanBTree+.


\subsection*{Question 9}
Quel est le taux de compression observé sur différents fichiers de texte ? Ce taux de compression dépend-il de la taille du fichier ? Y a-t-il d'autres paramètres qui influencent ce taux de compression ? Décrivez précisément les différents essais que vous aurez effectués avec vos programmes de compression et de décompression de texte. Quel est le temps de réponse moyen de votre programme de compression en fonction de la taille du fichier de texte ? Quel est le temps de réponse moyen de votre
programme de décompression en fonction de la taille du fichier comprimé ? \textbf{Les temps calcul observés sont-ils cohérents avec la complexité temporelle des algorithmes de compression/décompression ? Argumentez.} \\

Oui le taux de compression dépend de la taille du fichier, mais aussi de la fréquence à laquelle certain caractère reviennent. En effet pour chaque fichier compresser il y aura un header au fichier compresser avec l'arbre de Huffman sous forme binaire. Si le fichier est relativement petit, ce mode d'emploi va déjà prendre une certaine place, voir même plus que le fichier original. Par contre, si un fichier est très imposant, ce header sera insignifiant par rapport à la taille du fichier compressé.

De plus, plus un fichier contient des caractères avec beaucoup d'occurrences, plus l'arbre de Huffman sera petit et plus les nouvelles représentations en bit des caractères seront petites. La compression sera donc beaucoup plus efficace.

\subsection*{Question 10}
L'algorithme de compression par un codage de Huffman est normalement sans perte, c'est-à-dire que le fichier compressé contient toute l'information issue du fichier original. Une fois décompressé ce fichier doit donc être strictement identique à l'original. Votre implémentation garantit-elle cette propriété ? Pourquoi ? Quels sont les tests que vous avez effectués expérimentalement à ce propos ? \\

Notre implémentation garantit que tout l'information du fichier original sera restituée grâce à plusieurs choses : \\

\begin{itemize}
	\item Le nombre de caractères à écrire est contenu dans le fichier compressé. Dès lors, l'implémentation tient compte de cette donnée pour décompresser le fichier et ainsi récupérer la taille initiale du fichier avant compression.
	\item L'arbre de Huffman est contenu dans le fichier compressé sous forme binaire, et cette donnée permet de reconstruire exactement l'arbre original utilisé pour compresser le fichier. La traduction inverse de chaque séquence de bits du fichier compressé produira donc le caractère correct dont elle est issue.\\
\end{itemize}

Les tests que nous avons effectués pour en être certain sont basiques : il s'agit d'une simple succession de compression-décompression. Nous avons vérifié que les tailles des fichiers respectifs étaient égales et que leur contenu ne différait pas. \\

\section{Analyse de la complexité calculatoire}
\subsection{Complexité temporelle}

\section{Répartition du travail}

\begin{itemize}
\item Rédaction du rapport : Jonathan, Sundeep
\item Conception du programme : Boris, Tanguy, Alexandre, Mathieu
\end{itemize}

\section{Difficultés rencontrées}
Grâce aux questions de la séance intermédiaire, nous avons décortiqué le problème ce qui nous a permis de préparer correctement le partie implémentation. Nous n'avons donc pas rencontré de réel problème lors de la partie développement du programme.

\end{document}
